Fix a realized graph \(A\) in the support and write
\[
 I_0=I_0(A)=\sum_{i\in\mathcal V_n}\mathbf1\{N_i=0\}.
\tag{1}
\]
Whenever \(\mathcal M_{n,d,L}\ne\varnothing\), retain the graph primitives of any member law and replace only its response law by
\[
 f_i(a,x)=\varepsilon_iBx,
 \qquad
 \Pr(\varepsilon_i=1)=\Pr(\varepsilon_i=-1)=\frac12,
\tag{2}
\]
where the signs are independent across units and independent of the uniform latent types.  Then the pairs \((U_i,f_i)\) are independent and identically distributed, as required by \({\rm ass:unit\text{-}sampling}\).  For every \(a\in\{0,1\}\) and \(x\in[0,1]\),
\[
 |f_i(a,x)|\leq B,\qquad
 |\partial_x f_i(a,x)|=B,\qquad
 \partial_x^2f_i(a,x)=\partial_x^3f_i(a,x)=0.
\tag{3}
\]
Thus \({\rm ass:response\text{-}envelope}\) and \({\rm ass:third\text{-}derivative\text{-}bound}\) hold for every \(L\geq0\); the retained graph, assignment, rank, positivity, and overlap conditions are unchanged.  Hence (2) is a submodel of the maintained class.

For an isolate, \({\rm ass:anonymous\text{-}interference}\) gives
\[
 Y_i^{\mathrm{obs}}=f_i(Z_i,0)=0,
\tag{4}
\]
whereas its summand in \({\rm def:target}\) is
\[
 \frac{\partial_xf_i(1,1/2)+\partial_xf_i(0,1/2)}2
 =\varepsilon_iB.
\tag{5}
\]
Fix all nonisolate response functions.  Consider two supported sign configurations which assign, respectively, \(+1\) and \(-1\) to every isolate.  By (4), the two configurations induce the same conditional law of all observed data given \(A\), but by (5) their targets \(\theta_+\) and \(\theta_-\) satisfy
\[
 |\theta_+-\theta_-|=\frac{2B}{n}I_0.
\tag{6}
\]
Let \(m\) denote the common conditional expectation of any data-measurable estimator \(\widetilde\theta\) under these two configurations.  The triangle inequality gives
\[
 \max\{|m-\theta_+|,|m-\theta_-|\}
 \geq\frac{|\theta_+-\theta_-|}{2}
 =\frac BnI_0.
\tag{7}
\]
Both configurations have positive conditional-support probability under (2), independently of the graph.  Therefore every graph-only radius required to dominate the conditional bias for every response configuration in that support must satisfy
\[
 b(A)\geq\frac BnI_0(A).
\tag{8}
\]

Now give a genie all nonisolate signs.  Conditional on \(A\), the remaining isolate signs are independent Rademacher variables and are independent of the observed sample because their outcomes are identically zero by (4).  After subtracting the known nonisolate target contribution, estimation reduces exactly to the location problem
\[
 \frac BnS_{I_0},\qquad S_m=\sum_{k=1}^m\varepsilon_k.
\tag{9}
\]
The law of \(S_m\) is symmetric about zero, so zero is a median and minimizes expected absolute loss.  The conditional Bayes risk is therefore \((B/n)\mathbb E|S_{I_0}|\).  Independence and centering give
\[
 \mathbb ES_m^2=m,\qquad
 \mathbb ES_m^4=m+6\binom m2=3m^2-2m\leq3m^2.
\tag{10}
\]
Cauchy--Schwarz gives \(\mathbb E|S_m|\leq\sqrt m\).  Hölder's inequality gives
\[
 m=\mathbb ES_m^2
 \leq(\mathbb E|S_m|)^{2/3}(\mathbb ES_m^4)^{1/3}
 \leq(\mathbb E|S_m|)^{2/3}(3m^2)^{1/3},
\]
and hence, for \(m\geq1\),
\[
 \sqrt{m/3}\leq\mathbb E|S_m|\leq\sqrt m.
\tag{11}
\]
The conclusion is also valid at \(m=0\).  Substitution into (9) proves the two asserted conditional Bayes-risk bounds.  Averaging them over \(A\) yields the unconditional missing-sign fluctuation scale \(Bn^{-1}\mathbb E\sqrt{I_0(A)}\).  Since \(\sqrt m\leq m\) for every integer \(m\geq1\), this fluctuation is no larger than \(Bn^{-1}\mathbb EI_0(A)\).  The first argument concerns a graph-conditional, graph-only bias certificate, and the second is a genie-aided Bayes subproblem; neither step constructs an unconditionally honest interval or proves a universal expected-length lower bound for all such intervals.